\section{Related Work}
\label{sec:related_work}

We organize related work along three axes: (\S\ref{ssec:rw_paradigms}) the
LLM zero-shot reranking paradigms, with emphasis on setwise; (\S\ref{ssec:rw_extensions})
recent setwise / tournament extensions that share the ``more information per
LLM call'' goal but stay anchored to bounded local groups, including
\textbf{TourRank}; and (\S\ref{ssec:rw_long_context}) long-context attention
and position-bias work that constrains what is realistic at $N{=}50$.

\subsection{LLM Zero-Shot Reranking Paradigms}
\label{ssec:rw_paradigms}

LLM-based zero-shot reranking has converged on four prompting paradigms.
\textbf{Pointwise} methods score each query-document pair
independently~\citep{nogueira2020monot5,pradeep2021expandomonoduo}; they are
fast and easy to parallelize, but ignore inter-document context.
\textbf{Pairwise} methods compare two documents per
call~\citep{qin2024prp}; they extract the strongest local signal but pay
$O(n^2)$ comparisons in the all-pairs variant, mitigated by sorting-based
heapsort and bubblesort wrappers. \textbf{Listwise}
methods~\citep{sun2023rankgpt,ma2023zsl,zhang2025rankwithoutgpt,ren2025selfcalibrated}
feed a sliding window of candidates and ask the LLM to emit a permutation;
they are strong at moderate cost but suffer from output parsing fragility,
window-boundary effects, and well-documented position
bias~\citep{liu2024lostmiddle}. \textbf{Setwise}
methods~\citep{zhuang2024setwise} sit between pairwise and listwise: each
LLM call processes a small set of $c{+}1$ passages and emits a single label
identifying the most relevant one; the comparator is then embedded in
heapsort or bubblesort to produce a top-$k$ ranking. We adopt the setwise
paradigm and the public reference
implementation~\citep{ielab_llmrankers} as our baseline lineage.

The original setwise formulation~\citep{zhuang2024setwise} fixes
$c{+}1{=}4$, uses single-letter labels (\texttt{A}, \texttt{B}, \texttt{C},
\texttt{D}), and supports both \emph{generation} scoring (decode the winning
label) and \emph{likelihood} scoring (compare logits of candidate label
tokens, on encoder-decoder backbones). MaxContext extends this lineage in
three respects: (i) the prompt feeds the entire live rerank pool rather than
a small window; (ii) labels are numeric $1$--$N$ to support
$N{\,>\,}23$ under strict parser invariants; and (iii) the comparator can
emit either one (\textsc{TopDown}, \textsc{BottomUp}) or two
(\textsc{DualEnd}) labels per call. We use generation scoring throughout
because joint best-worst elicitation requires the model to emit two distinct
labels in a single response, which the likelihood scoring path cannot
represent without collapsing to a best-only proxy.

\subsection{Setwise and Tournament Extensions}
\label{ssec:rw_extensions}

Several recent papers share MaxContext's high-level goal of extracting more
ranking evidence per LLM call, but each is anchored to bounded local groups
of $c{+}1{\le}20$.

\paragraph{Joint best-worst elicitation.}
Building on the public setwise reference
implementation~\citep{ielab_llmrankers}, we implement a \textsc{DualEnd}
comparator that asks the LLM to identify the most \emph{and} least relevant
document in a single prompt, doubling the explicit information per call from
one label to a partition of the $c{+}1$ candidates into two extremes plus
the rest. The dual-output comparator naturally pairs with cocktail-shaker
sort or double-ended selection sort. In our completed baseline experiments
on TREC DL19/20 across nine LLMs, DualEnd is the strongest small-window
family but pays $5.6{-}8.9\times$ the wall-clock of TopDown-Heap. We treat
small-window DualEnd as a \emph{baseline} for this paper rather than a
contribution: MaxContext-DualEnd inherits the joint best-worst comparator
and applies it to the entire live pool instead of a $c{+}1{=}4$ window.

\paragraph{Tournament prompting: TourRank.}
\textbf{TourRank}~\citep{chen2025tourrank} runs setwise comparisons inside a
multi-round tournament bracket. Documents are seeded into groups using BM25
ranks, each round uses an LLM to advance a fraction of each group's
candidates to the next round, and a final ranking is produced by accumulated
points across $r$ tournaments. The information advantage over plain setwise
heapsort is structural: a single comparator output provides broader ranking
evidence through advancement, elimination, and accumulated points across
rounds, and grouping seeded by initial rank reduces the chance that all
strong candidates are concentrated in one group. TourRank is highly
parallelizable --- groups within a round are independent, and tournaments
are independent of each other --- and reports state-of-the-art quality on
TREC DL and BEIR at moderate call budgets. Recent
work~\citep{blitzrank2026} generalizes the tournament intuition into
principled tournament \emph{graphs} that explicitly model all preferences
extractable from a $k$-way comparison, rather than only winner identity.
MaxContext is complementary to both threads: it does not change the bracket
/ graph topology, but instead changes the \emph{size} of the live group
inside which a setwise comparator operates, pushing it from $c{+}1{\le}20$
all the way to $N{=}50$.

\paragraph{Other setwise efficiency improvements.}
\textbf{Setwise Insertion}~\citep{podolak2025setwiseinsertion} warm-starts the
sort using the BM25 ranking as a prior, reducing query latency by
$\sim31\%$ and the number of inferences by $\sim23\%$ without changing the
comparator. This is orthogonal to MaxContext: an insertion-style warm start
could be applied within MaxContext's pool-shrinking loop, though we leave
that to future work. \textbf{Rank-R1}~\citep{zhuang2025rankr1} extends the
setwise comparator with a reasoning chain learned via reinforcement learning
and uses $20$ candidates per call --- a known-feasible precedent for
moderately large setwise windows on Qwen-class models. Rank-R1 is
complementary: it changes the \emph{quality} of each comparison via training,
while MaxContext changes the \emph{scale} of each comparison via
prompt-engineering at zero-shot. Combining a reasoning-trained comparator
with whole-pool prompting is again future work.
\textbf{REALM}~\citep{wang2025realm} models per-document relevance as
Gaussian distributions and uses recursive Bayesian updates to reduce LLM
queries by 23--84\%. MaxContext takes the opposite extreme of REALM's design
space: rather than adapting the sort to the comparator's noise, it makes the
comparator's window so large that the sort itself collapses to a few rounds
of ``pick-and-place from the live pool.''

\paragraph{Other paradigm-level alternatives.}
\textbf{Attention-based reranking}~\citep{chen2025icr} skips generation
entirely and uses attention scores from the LLM's forward pass as relevance
weights, giving order-of-magnitude latency reductions at some quality cost;
it is a different paradigm from setwise prompting.
\textbf{BlockRank}~\citep{gupta2025blockrank} exploits inter-document block
sparsity in attention to scale in-context ranking to long candidate lists
within a single forward pass, emphasising raw throughput.
MaxContext differs from both in that it remains within the setwise prompting
paradigm and pays the cost of repeated whole-pool generation calls in
exchange for the ability to elicit best, worst, or both extremes from a
long-context comparator. Reasoning-enhanced rerankers such as
\textbf{Rank1}~\citep{rank1_2025} use test-time compute to improve a
pointwise comparator; the techniques are stackable but address a different
axis (per-call quality, not per-query call count).

\subsection{Long-Context Attention, Position Bias, and Order Robustness}
\label{ssec:rw_long_context}

Whole-pool prompting at $N{=}50$ pushes a $\sim25$K-token prompt through the
LLM, which exposes MaxContext to two well-studied risks.

\paragraph{Lost-in-the-middle.}
\citet{liu2024lostmiddle} show that LLMs disproportionately attend to
information at the beginning and end of long contexts, with a documented
``U-shaped'' performance curve. \citet{hutter2025positionalrag} extend the
phenomenon to retrieval-augmented generation, finding that bias depends on
document type and model architecture. \citet{tang2024foundmiddle} propose
\emph{permutation self-consistency} for listwise reranking: shuffle the
input list multiple times, run the LLM on each permutation, and aggregate
the rankings. This mitigates position bias at the cost of multiplying
inference. MaxContext does \emph{not} apply permutation self-consistency by
default because doing so would multiply the call count we are trying to
reduce; instead, we treat order robustness as a launch gate, running a
controlled-ordering pilot (BM25-forward, BM25-inverse, random) at
$N{=}50$ and reporting the maximum pairwise NDCG@10 difference as evidence
for or against order stability at scale.

\paragraph{LLM comparator inconsistency.}
\citet{zeng2024llmrankfusion} document two intrinsic inconsistencies in
LLM-based comparison: order inconsistency (swapping inputs gives different
outputs) and transitivity inconsistency (preference cycles). Sorting-based
setwise approaches assume both. MaxContext does not eliminate these
inconsistencies, but it reduces \emph{exposure} to them: with only
$\lfloor N/2 \rfloor$ or $N{-}2$ comparator calls per query, fewer
opportunities for inconsistent outputs accumulate compared to small-window
sorts that make several hundred calls. Order sensitivity is diagnosed by the
controlled-ordering pilot in \S\ref{sec:experimental_setup}; format
reliability is diagnosed separately by parse-fallback rates at $N{=}50$.

\paragraph{Sorting with noisy comparators and efficiency reporting.}
\citet{sato2026sortingsurvey} surveys sorting algorithms specifically for
LLM-based comparators, noting that classical complexity bounds may not be
optimal when comparisons are expensive, batchable, and noisy. Recent work
also argues that LLM-reranker efficiency cannot be summarized by call count
alone. \citet{peng2025flopsreranking} introduce hardware-independent
ranking-per-PetaFLOP and queries-per-PetaFLOP metrics, recognising that a
single LLM call with $25$K prompt tokens is fundamentally different from
$25$ small-window calls with $1$K tokens each. We adopt the spirit of this
recommendation: throughout the experimental sections we report all five cost
axes (LLM calls, prompt tokens, completion tokens, total tokens, wall-clock)
on every method, and frame our claim narrowly as potential Pareto-frontier
movement on the comparisons-axis and wall-clock-axis rather than on
per-token efficiency.

\subsection{Positioning of MaxContext}

Within this landscape, MaxContext is the first setwise reranker we are aware
of that prompts over the entire live rerank pool on every call (up to
$N{=}50$ documents) while retaining a controlled comparison of comparator
output type via the three-variant ablation. It is complementary to
TourRank-style tournament structure (which changes \emph{how} comparator
outputs are aggregated), to Setwise Insertion (which changes \emph{where}
the sort starts), and to Rank-R1 (which changes \emph{what} the comparator
internally computes). Unlike permutation self-consistency methods, MaxContext
treats order robustness as an empirical constraint rather than as a
mitigation strategy: we evaluate controlled input orders rather than
multiplying inference by default. The natural follow-up question --- whether
stacking MaxContext with a tournament structure or a reasoning-trained
comparator yields further improvements --- is left to future work.
